\section{Economic Study}

\subsection{Introduction}
spiegazione della questione di 0 tassi di interesse per tipo di progetto (prestito tax free di UE)

\subsection{CAPEX}
capex con fonti + ipotesi decrescita capex 

\subsection{OPEX}
opex con fonti + ipotesi decrescita opex

\subsection{Revenues}
ricavi con fonti + ipotesi crescita/decrescita ricavi e calcolo PUN etc.. 

\subsection{Economic Results: base case}

\subsubsectionn{Sensitivity Analysis}

\subsection{Economic Results: alternative scenarios}
dato che buco finanziario, i primi 15 anni facciamo solo reattore (full electric) per mettere da parte un po' di profitto
quando la possibilità di decarbonizzare il settore dei trasporti è più matura, dopo 15 anni si costruisce l'impianto secondario 
con HTSE, metanatore e turbina a gas. questo permette di guadagnare di più prima e investire più tardi sui componenti secondari.

risultato è uno scenario ibrido con NPV positivo, con 5 anni in meno di operation per ottimizzare il ricambio di HTSE.



